% Section III: Mathematical toolkit - just-in-time math, not front-loaded
% Purpose: Concise conventions and prerequisites
% Full mathematical details distributed to sections where they're used
%
% Organization principle: Introduce concepts just before they're needed
% - π₃(G) appears in Section V with level quantization
% - Holonomy and flat connections appear in Section V with classical solutions
% - Symplectic quotient appears in Section VI with quantization
% - Bordisms appear in Section VII with functorial TQFT
%
% Status: EXTRACT AND REORGANIZE from topology_review.tex
% Sources:
% - tex_docs/topology_review.tex (comprehensive math, needs to be broken up)
% - Keep only: basic conventions, notation, minimal background
% - Move detailed treatments to sections where they're applied
%
% This section should be SHORT: ~5 pages maximum
% More like "mathematical conventions and prerequisites" than a mini-textbook

\section{Mathematical Conventions and Prerequisites}

% TODO: Extract minimal necessary background from topology_review.tex
% Move most content to later sections where it's used

\subsection{Notation and Conventions}

% TODO: Manifold notation, group conventions, what reader should already know

\subsection{Differential Forms: Brief Review}

% TODO: Very brief - wedge product, exterior derivative, Hodge star
% Reader should already know this from E&M
% Full treatment in Section IV when building CS action

\subsection{Lie Groups and Algebras}

% TODO: G, 𝔤, adjoint action, structure constants
% Just enough to read CS action
% Deeper treatment when needed in Section V

\subsection{Gauge Fields as Connections}

% TODO: A as 𝔤-valued 1-form, curvature F, gauge transformations
% Bridge from physics notation to geometric language
% This is the most important subsection

% Note: Homotopy groups, holonomy, moduli spaces, symplectic geometry, bordisms
% all appear later when they're actually used
